% !TeX encoding = UTF-8
% !TeX root = ../main.tex

\section{信号分析与特征提取}

\begin{frame}{信号分析：时域、频域与时频表示}
  \begin{columns}[T]
    \column{0.48\textwidth}
    \centering
    \includegraphics[width=0.80\textwidth]{waveform.png}\\[-6pt]
    {\footnotesize (a) 时域波形}

    \vspace{4pt}
    \includegraphics[width=0.80\textwidth]{fft_spectrum.png}\\[-6pt]
    {\footnotesize (b) FFT 幅度谱（0--4000 Hz）}

    \column{0.48\textwidth}
    \centering
    \includegraphics[width=0.80\textwidth]{spectrogram.png}\\[-6pt]
    {\footnotesize (c) STFT 语谱图}

    \medskip
    \raggedright\small
    语音是非平稳信号，单一 FFT 丢失频率随时间变化的信息。STFT 通过分帧加窗实现时频联合分析：
    \[
    \operatorname{STFT}\{x[n]\}(m,\omega) = \sum_{n=-\infty}^{\infty} x[n] w[n-mR] e^{-j\omega n}
    \]
    {\footnotesize 25 ms Hamming 窗，10 ms 帧移}
  \end{columns}
\end{frame}

\begin{frame}{MFCC 特征提取流程}
  \small
  \begin{columns}[T]
    \column{0.50\textwidth}
    \textbf{梅尔频率倒谱系数（MFCC）}
    \begin{enumerate}
      \setlength{\itemsep}{1pt}
      \item \textbf{预加重}：$y[n] = x[n] - 0.97 x[n-1]$（FIR 高通）
      \item \textbf{分帧加窗}：25 ms Hamming 窗，10 ms 帧移
      \item \textbf{FFT $\to$ 功率谱 $\to$ Mel 滤波}：26 个三角滤波器（低频密、高频疏）
      \item \textbf{取对数 + DCT}：压缩动态范围，去相关，取前 13 维
      \item \textbf{动态特征}：$\Delta$ 和 $\Delta^2$（一阶/二阶差分）$\to$ 39 维
    \end{enumerate}

    \column{0.46\textwidth}
    \centering
    \includegraphics[width=0.80\textwidth]{mfcc.png}\\[-4pt]
    {\footnotesize 13 维 MFCC（传统方法）}

    \vspace{8pt}
    \includegraphics[width=0.80\textwidth]{mfcc_features.png}\\[-4pt]
    {\footnotesize 39 维 MFCC + $\Delta$ + $\Delta^2$（深度方法）}
  \end{columns}
\end{frame}
